\section*{Chapter \ref{ch:b1:heat-budget} --- The Earth's Heat Budget}

\begin{solution}{exo:b1:heat-budget:1}
$q = k\,\mathrm dT/\mathrm dz = 2.5\times\qty{0.022}{K.m^{-1}} =
\qty{0.055}{W.m^{-2}} = \qty{55}{mW.m^{-2}}$. Through \qty{1}{km^2} $=
\qty{e6}{m^2}$: \qty{5.5e4}{W}, that is \qty{55}{kW}.
\end{solution}

\begin{solution}{exo:b1:heat-budget:2}
Continents $\num{2.073e14}\times0.0709 = \qty{1.470e13}{W} = \qty{14.70}{TW}$;
oceans $\num{3.028e14}\times0.1054 = \qty{31.92}{TW}$; total \qty{46.61}{TW}, and
the mean \omterm{def:b1:heat-budget:heat-flow}{heat flow} $\num{4.661e13}/\num{5.101e14} = \qty{91.4}{mW.m^{-2}}$ (the
published 46.7 and 91.6 differ only by the rounding of the two means). The
fraction of the solar input is $46.6/173000 = \num{2.7e-4}$, or
\qty{0.027}{\percent}.
\end{solution}

\begin{solution}{exo:b1:heat-budget:3}
Uranium: $0.99274\times\num{9.513e-5} + 0.00720\times\num{5.685e-4} =
\num{9.444e-5} + \num{4.09e-6} = \qty{9.85e-5}{W.kg^{-1}}$; the rare
\ce{^{235}U} supplies \qty{4}{\percent} of it. Potassium: the mass fraction of
\ce{^{40}K} is $\num{1.17e-4}\times39.96/39.10 = \num{1.196e-4}$, so
$\num{1.196e-4}\times\num{2.847e-5} = \qty{3.40e-9}{W.kg^{-1}}$.
\end{solution}

\begin{solution}{exo:b1:heat-budget:4}
$\qty{1}{\micro cal.cm^{-2}.s^{-1}} = \num{4.184e-6}\,\unit{J.s^{-1}}/\qty{e-4}{m^2}
= \qty{41.84}{mW.m^{-2}}$: eastern United States \qty{33.5}{mW.m^{-2}}, Basin and
Range \qty{58.6}{mW.m^{-2}}, Sierra Nevada \qty{16.7}{mW.m^{-2}}. With $A_0D$ in
\unit{mW.m^{-2}} when $A_0$ is in \unit{\micro W.m^{-3}} and $D$ in \unit{km}:
$33.5 + 3.0\times7.5 = \qty{56.0}{mW.m^{-2}}$, $58.6 + 3.0\times9.4 =
\qty{86.8}{mW.m^{-2}}$, $16.7 + 3.0\times10 = \qty{46.7}{mW.m^{-2}}$. The same
granite gives \omterm{def:b1:heat-budget:heat-flow}{heat flows} almost twice as high in the Basin and Range as in the
Sierra Nevada: the difference is in the \omterm{def:b1:layered-earth:layers}{mantle} below.
\end{solution}

\begin{solution}{exo:b1:heat-budget:5}
$T(10\,\text{km}) = 10 + \frac{0.060\times10^4}{2.5} - \frac{\num{2e-6}\times10^8}{2\times2.5}
= 10 + 240 - 40 = \qty{210}{\degreeCelsius}$. Below, $q_m = 60 - 2.0\times10 =
\qty{40}{mW.m^{-2}}$, a gradient of \qty{16}{K.km^{-1}}, and
$T(35\,\text{km}) = 210 + 16\times25 = \qty{610}{\degreeCelsius}$.
\end{solution}

\begin{solution}{exo:b1:heat-budget:6}
Shale: gradient $12.0/0.300 = \qty{40}{K.km^{-1}}$, $q = 1.5\times0.040 =
\qty{60}{mW.m^{-2}}$. Granite: $10.0/0.500 = \qty{20}{K.km^{-1}}$,
$q = 3.0\times0.020 = \qty{60}{mW.m^{-2}}$. $R(300) = 300/1.5 =
\qty{200}{m^2.K.W^{-1}}$, $R(800) = 200 + 500/3.0 = \qty{366.7}{m^2.K.W^{-1}}$.
The three points $(0, 12.0)$, $(200, 24.0)$, $(366.7, 34.0)$ are aligned, with
slope $22.0/366.7 = \qty{0.0600}{W.m^{-2}}$: the \omterm{def:b1:heat-budget:heat-flow}{heat flow} is the same,
\qty{60}{mW.m^{-2}}, in both rocks; only the gradient changes, in inverse
proportion to the conductivity.
\end{solution}

\begin{solution}{exo:b1:heat-budget:7}
$\tau = L^2/\kappa$: \qty{1}{m}, \qty{e6}{s} or \num{12}~days; \qty{10}{m},
\qty{e8}{s} or \num{3.2}~years; \qty{100}{km}, \qty{e16}{s} or
\qty{317}{Myr}; \qty{2900}{km}, \qty{8.4e18}{s} or \qty{267}{Gyr}. At
\qty{10}{m} the daily cycle is not felt at all, and the yearly cycle, shorter
than the three years needed to cross \qty{10}{m}, only very weakly: only changes
lasting years or more come through. The last time is $267/4.54 = 59$ times the
age of the Earth.
\end{solution}

\begin{solution}{exo:b1:heat-budget:8}
(a) $1646 - 0.54\times50 = \qty{1619}{K}$, that is \qty{1346}{\degreeCelsius},
the published \qty{1350}{\degreeCelsius} rounded. (b) Mean gradient
$(0.54 + 0.36)/2 = \qty{0.45}{K.km^{-1}}$, so $1646 + 0.45\times360 =
\qty{1808}{K}$, \qty{9}{K} (\qty{0.5}{\percent}) above \qty{1799}{K}: the
gradient does not decrease linearly, but the estimate is close.
\end{solution}

\begin{solution}{exo:b1:heat-budget:9}
(a) $G^2 = T_0^2/(\pi\kappa t)$ gives $t = T_0^2/(\pi\kappa G^2)$.
(b) $G = 0.02\times\tfrac59\,\unit{K}/\qty{0.3048}{m} = \qty{0.03645}{K.m^{-1}} =
\qty{36.45}{K.km^{-1}}$; $t = 3900^2/(\pi\times\num{1.2e-6}\times0.03645^2) =
\num{1.521e7}/\num{5.009e-9} = \qty{3.04e15}{s}$, that is \qty{96}{Myr},
Kelvin's figure. (c) $t = 1300^2/(\pi\times\num{e-6}\times0.036^2) =
\qty{4.15e14}{s} = \qty{13}{Myr}$: a more realistic initial temperature makes
the conducting Earth younger still.
\end{solution}

\begin{solution}{exo:b1:heat-budget:10}
(a) After erosion, the new surface is the old depth $e$, where the heat
production was $A_0' = A_0e^{-e/D}$ and the \omterm{def:b1:heat-budget:heat-flow}{heat flow} $q(e) = q_m +
A_0De^{-e/D} = q_m + A_0'D$; measured from the new surface, the production
still decreases as $A_0'e^{-z'/D}$. The pluton moves down along the same line.
(b) The layer is now $h - e$ thick with the same $A$, so $q_0' = q_m + A(h - e)$:
on a line of slope $h - e$ through $q_m$. (c) Only the exponential model: in the
uniform-layer model, plutons eroded by different depths would lie on lines of
different slopes, and the points of one province would scatter. That a single
line is observed favours a production decreasing roughly exponentially with
depth.
\end{solution}

\begin{solution}{exo:b1:heat-budget:11}
(a) $\tau_P = RL/(3\kappa)$: for \qty{50}{km}, $\num{6.4e6}\times\num{5e4}/\num{3e-6}
= \qty{1.07e17}{s} = \qty{3.4}{Gyr}$; for \qty{100}{km}, \qty{6.8}{Gyr}.
(b) $q = kT_i/L = 3\times1300/\num{5e4} = \qty{78}{mW.m^{-2}}$ at the start,
and $78\,e^{-2/3.38} = \qty{43}{mW.m^{-2}}$ after \qty{2}{Gyr}. (c)
$3\times1300/\sqrt{\pi\times\num{e-6}\times\num{6.31e16}} = 3900/\num{4.45e5} =
\qty{8.8}{mW.m^{-2}}$: the well-stirred Earth still loses five times more heat
after two billion years. Its whole interior feeds the surface, where only an
outer shell does in a conducting Earth.
\end{solution}

\begin{solution}{exo:b1:heat-budget:12}
(a) $\sqrt{\pi\kappa\times\qty{1}{Myr}} = \sqrt{\pi\times\num{e-6}\times\num{3.156e13}}
= \qty{9957}{m}$, so $C = 3\times1336/9957 = \qty{0.403}{W.m^{-2}}$ for
$t$ in Myr, \qty{403}{mW.m^{-2}.Myr^{1/2}}: \qty{82}{\percent} of the fitted
490. The model needs a larger $k\,\Delta T$ (about \qty{4900}{W.m^{-1}}): a
higher conductivity of the \omterm{def:b1:layered-earth:layers}{mantle}, or a hotter interior, as along the adiabat.
(b) $490/\sqrt1 = 490$, $490/\sqrt{10} = 155$ and $490/\sqrt{100} =
\qty{49}{mW.m^{-2}}$. (c) $z = 2.326\sqrt{\num{e-6}\times\num{3.156e15}} =
2.326\times\qty{56180}{m} = \qty{131}{km}$.
\end{solution}

\begin{solution}{pb:b1:heat-budget:1}
\textbf{1.} Gradient $= (T_{500} - T_{100})/\qty{0.4}{km}$: A \qty{15.25}{K.km^{-1}},
B \num{19.50}, C \num{19.50}, D \num{25.50}, E \qty{28.25}{K.km^{-1}}.

\textbf{2.} Drilling disturbs the temperature of the rock (the drilling fluid
cools or warms the walls), and the disturbance takes weeks to months to die away.
The top tens of metres follow the seasonal cycle and recent changes of surface
temperature.

\textbf{3.} $q_0 = k\times$gradient: A $2.9\times15.25 = \qty{44.225}{mW.m^{-2}}$,
B \qty{50.70}{mW.m^{-2}}, C \qty{60.45}{mW.m^{-2}}, D \qty{68.85}{mW.m^{-2}},
E \qty{84.75}{mW.m^{-2}}.

\textbf{4.} A to D are below the continental mean, E above it by
\qty{13.9}{mW.m^{-2}}, or \qty{20}{\percent}.

\textbf{5.} The \omterm{def:b1:heat-budget:heat-flow}{heat flow} is conductivity times gradient: C conducts better (3.1
against 2.6), so the same gradient carries \qty{19}{\percent} more heat.

\textbf{6.} The isotherms follow the relief only partly: under a summit they
rise less than the surface, so they are farther apart beneath it, and the
measured gradient is lowered (raised under a valley). A correction for relief is
needed before comparing the plutons.

\textbf{7.} $\rho C_U H_U$: A $2650\times\num{2.0e-6}\times\num{9.85e-5} =
\qty{0.522}{\micro W.m^{-3}}$, B 0.914, C 1.305, D 1.827, E 2.610.

\textbf{8.} Thorium: A 0.558, B 0.976, C 1.255, D 1.742, E 2.439.

\textbf{9.} Potassium: A $2650\times0.030\times\num{3.40e-9} = 0.270$, B 0.315,
C 0.360, D 0.378, E 0.405.

\textbf{10.} $A_0$ = 1.35, 2.20, 2.92, 3.95 and \qty{5.46}{\micro W.m^{-3}}.
Potassium contributes least, from \qty{20}{\percent} of the total in A down to
\qty{7}{\percent} in E; uranium and thorium share the rest almost equally.

\textbf{11.} Yes: the five points rise steadily, within about
\qty{1.5}{mW.m^{-2}} of a straight line.

\textbf{12.} $\bar A_0 = 15.88/5 = 3.176$, $\bar q_0 = 308.975/5 = 61.795$;
$\sum(A_0 - \bar A_0)^2 = 10.1681$; $\sum(A_0 - \bar A_0)(q_0 - \bar q_0) =
101.146$.

\textbf{13.} $D = 101.146/10.1681 = 9.95$ (in \unit{mW.m^{-2}} per
\unit{\micro W.m^{-3}}, that is in km): $D = \qty{9.95}{km}$, and $q_m = 61.795 -
9.95\times3.176 = \qty{30.20}{mW.m^{-2}}$.

\textbf{14.} The slope, about \qty{10}{km}, is that of the three provinces; the
reduced \omterm{def:b1:heat-budget:heat-flow}{heat flow}, \qty{30}{mW.m^{-2}}, is close to the eastern United States'
\qty{33}{mW.m^{-2}}: a quiet province with a cool \omterm{def:b1:layered-earth:layers}{mantle} below, unlike the Basin
and Range.

\textbf{15.} $q_m/k = \qty{12}{K.km^{-1}}$ and $A_0D^2/k = \num{5.46e-6}\times\num{e8}/2.5
= \qty{218.4}{K}$: $T(z) = 10 + 12z + 218.4\,(1 - e^{-z/10})$, $z$ in km, $T$ in
\unit{\degreeCelsius}.

\textbf{16.} $q(z) = 30 + 54.6\,e^{-z/10}$: \qty{50.1}{mW.m^{-2}} at \qty{10}{km},
\num{37.4} at \qty{20}{km} and \qty{31.7}{mW.m^{-2}} at \qty{35}{km}.

\textbf{17.} $30/84.6 = \qty{35}{\percent}$: nearly two thirds of the \omterm{def:b1:heat-budget:heat-flow}{heat flow}
over E is produced in its radiogenic layer.

\textbf{18.} $T(10) = 10 + 120 + 218.4\times0.632 = \qty{268}{\degreeCelsius}$.

\textbf{19.} $T(20) = 10 + 240 + 218.4\times0.865 = \qty{439}{\degreeCelsius}$.

\textbf{20.} Under A, $A_0D^2/k = 1.35\times100/2.5 = \qty{54.0}{K}$:
$T(35) = 10 + 420 + 54.0\times0.970 = \qty{482}{\degreeCelsius}$.

\textbf{21.} By $A_0D^2/k$: \qty{54}{K} under A and \qty{218}{K} under E. The deep
\omterm{def:b1:heat-budget:heat-flow}{heat flow} is the same; E's \omterm{def:b1:layered-earth:layers}{crust} is four times richer in heat sources, and the
heat they produce must be conducted up through the same rock, which needs a
steeper gradient all the way.

\textbf{22.} $1300 = 10 + 218.4 + 12z$ gives $z = \qty{89.3}{km}$ under E;
$1300 = 10 + 54.0 + 12z$ gives \qty{103.0}{km} under A.

\textbf{23.} Both lie within the range of continental \omterm{def:b1:layered-earth:lithosphere}{lithosphere}, from a
hundred kilometres or less under young regions to about \qty{250}{km} under
cratons, towards its thin side: with a reduced \omterm{def:b1:heat-budget:heat-flow}{heat flow} of \qty{30}{mW.m^{-2}} and this
conductivity, the conductive \omterm{def:b1:heat-budget:geotherm}{geotherm} reaches \omterm{def:b1:layered-earth:layers}{mantle} temperatures at about
\qty{100}{km}; under the more radioactive pluton, \qty{14}{km} sooner.

\textbf{24.} $T(35) = 10 + 420 + 218.4\times(1 - e^{-3.5}) = 430 + 211.8$: the
\omterm{def:b1:layered-earth:discontinuity}{Moho} beneath pluton E is at
\textbf{\qty[text-series-to-math=true]{642}{\degreeCelsius}}.
\end{solution}
